% Part: many-valued-logic
% Chapter: syntax-and-semantics
% Section: sublogics

\documentclass[../../../include/open-logic-section]{subfiles}

\begin{document}

\olfileid{mvl}{syn}{sub}

\olsection{多值逻辑作为~$\LogCL$ 的子逻辑}

通常的多值逻辑都是使用矩阵来定义的，在这些矩阵中，当自变元取自 $\{\True, \False\}$ 时，真值函数的值与经典真值函数一致。具体来说，在这些逻辑中，若 $x \in \{\True, \False\}$，则 $\tf{\lnot}[\Log L](x) =
\tf{\lnot}[\LogCL](x)$，并且对于 $\land$、$\lor$、$\lif$ 中的任一 $\star$，若 $x, y \in \{\True, \False\}$，则 $\tf{\star}[\Log L](x,y) =
\tf{\star}[\LogCL](x,y)$。换言之，$\lnot$、$\land$、$\lor$、$\lif$ 的真值函数限制在 $\{\True,\False\}$ 上时，即为经典真值函数。

\begin{prop}\ollabel{prop:mvl-cl}
  假设多值逻辑~$\Log L$ 的语言包含联结词
  $\lnot$、$\land$、$\lor$、$\lif$，$\True, \False \in
  V$，且其真值函数满足：
  \begin{enumerate}
    \item\ollabel{prop:not} $\tf{\lnot}[\Log L](x) = \tf{\lnot}[\LogCL](x)$（若 $x =
    \True$ 或 $x = \False$）；
    \item\ollabel{prop:land} $\tf{\land}[\Log L](x,y) = \tf{\land}[\LogCL](x,y)$，
    \item\ollabel{prop:lor} $\tf{\lor}[\Log L](x,y) = \tf{\lor}[\LogCL](x,y)$，
    \item\ollabel{prop:lif} $\tf{\lif}[\Log L](x,y) = \tf{\lif}[\LogCL](x,y)$，
    （若 $x, y \in \{\True, \False\}$）。
  \end{enumerate}
  那么，对于任何到~$V$ 中的赋值 $\pAssign v$，若 $\pAssign
  v(p) \in \{\True,\False\}$，便有 $\pValue v[\Log L](!A) = \pValue
  v[\LogCL](!A)$。
\end{prop}

\begin{proof}
  对~$!A$ 施归纳。
  \begin{enumerate}
  \item 若 $!A \ident p$ 为原子公式，则 $\pValue v[\Log L](!A) =
  \pAssign v(p) = \pValue
  v[\LogCL](!A)$。
  \item 若 $!A \ident \lnot B$，则
  \begin{align*}
    \pValue v[\Log L](!A) & = \tf{\lnot}[\Log L](\pValue v[\Log L](!B)) & &\text{根据 \olref[val]{defn:pValue}}\\
    & = \tf{\lnot}[\Log L](\pValue v[\LogCL](!B)) && \text{由归纳假设}\\
    & = \tf{\lnot}[\LogCL](\pValue v[\LogCL](!B))
&&\text{由假设 \olref{prop:not}，}\\
&&&\text{因为 $\pValue v[\LogCL](!B) \in \{\True, \False\}$，}\\
& = \pValue v[\LogCL](!A) &&\text{根据 \olref[val]{defn:pValue}}.
\end{align*}
\item 若 $!A \ident (!B \land !C)$，则
\begin{align*}
  \pValue v[\Log L](!A) & = \tf{\land}[\Log L](\pValue v[\Log L](!B), \pValue v[\Log L](!C)) & &\text{根据 \olref[val]{defn:pValue}}\\
  & = \tf{\land}[\Log L](\pValue v[\LogCL](!B),\pValue v[\LogCL](!C)) && \text{由归纳假设}\\
  & = \tf{\land}[\LogCL](\pValue v[\LogCL](!B),\pValue v[\LogCL](!C))
&&\text{由假设 \olref{prop:land}，}\\
&&&\text{因为 $\pValue v[\LogCL](!B),\pValue v[\LogCL](!C) \in \{\True, \False\}$，}\\
& = \pValue v[\LogCL](!A) &&\text{根据 \olref[val]{defn:pValue}}.
\end{align*}
\end{enumerate}
$!A \ident (!B \lor !C)$ 和 $!A \ident (!B \lif !C)$ 的情形类似。
\end{proof}

\begin{cor}
  若多值逻辑满足 \olref{prop:mvl-cl} 的条件，且 $\True \in V^+$ 而 $\False \notin V^+$，则 ${\Entails[\Log L]} \subseteq {\Entails[\LogCL]}$，即若 $\Gamma
  \Entails[\Log L] !B$，则 $\Gamma \Entails[\LogCL] !B$。特别地，$\Log L$ 的每个重言式也是经典重言式。
\end{cor}

\begin{proof}
  我们证明其逆否命题。设 $\Gamma \Entails/[\LogCL] !B$。
  则存在赋值~$\pAssign v\colon \PVar \to
  \{\True, \False\}$，使得对所有 $!A \in \Gamma$ 有 $\pValue v[\LogCL](!A) = \True$，且 $\pValue v[\LogCL](!B) = \False$。由于 $\True,
  \False \in V$，赋值~$\pAssign v$ 也是~$\Log L$ 的赋值。
  由 \olref{prop:mvl-cl}，对所有 $!A \in \Gamma$ 有 $\pValue v[\Log L](!A) =
  \True$，且 $\pValue v[\Log L](!B) = \False$。
  由于 $\True \in V^+$ 且 $\False \notin V^+$，这意味着 $\pAssign v
  \Entails[\Log L] \Gamma$ 且 $\pAssign v \Entails/[\Log L] !B$，
  即 $\Gamma \Entails/[\Log L] !B$。
\end{proof}
\end{document}